\documentclass[11pt]{article}

\usepackage[margin=1in]{geometry}
\usepackage{amsmath,amssymb,amsthm,mathtools}
\usepackage{microtype}
\usepackage{enumitem}
\usepackage{xcolor}
\usepackage{xr-hyper}
\usepackage{hyperref}
\usepackage[nameinlink,noabbrev]{cleveref}

\newcommand{\Ztwo}{\mathbb Z_2}

\externaldocument[red-]{source_internal_control_factorization_note}
\externaldocument[stack-]{odd_characterization_transport_stack_note}
\externaldocument[splitlift-]{split_controlled_completion_lift_note}

\hypersetup{
  colorlinks=true,
  linkcolor=blue!60!black,
  citecolor=blue!60!black,
  urlcolor=blue!60!black
}

\theoremstyle{definition}
\newtheorem{definition}{Definition}[section]
\theoremstyle{plain}
\newtheorem{theorem}[definition]{Theorem}
\newtheorem{proposition}[definition]{Proposition}
\newtheorem{assumption}[definition]{Assumption}
\newtheorem{corollary}[definition]{Corollary}
\theoremstyle{remark}
\newtheorem{remark}[definition]{Remark}

\title{Delayed Odd-Control and Next-Receiver Readout Note}
\author{}
\date{}

\begin{document}
\maketitle

Fix an accepted tick \(k\) and the next accepted receiver step \(k+1\).
This note isolates the delayed odd/control datum remaining after current accepted-record reduction,
defines the corresponding delayed quotient and selected line, constructs the delayed realization map
to the next-step reduced welded LoG line, and packages selected-line delayed response sufficiency
with an explicit imported nondegeneracy clause.

\section{Current Reduced Receiver Closure}

\begin{proposition}[Current reduced receiver closure with delayed odd residue]
\label{prop:delay-current-closure}
Fix an accepted welded branch and an accepted tick \(k\).
Assume the accepted-record reduction theorem \cref*{red-thm:accepted-record-reduction}.
Then the current theorem-visible receiver object is closed under accepted-record reduction:
\begin{enumerate}[leftmargin=2em,label=(\roman*)]
\item equality of accepted-record history up to the current accepted tick fixes every current
  theorem-visible reduced receiver datum;
\item any remaining realization dependence can live only outside the current theorem-visible reduced
  receiver object, namely in the off-diagonal complement or in the leakage/slack branch;
\item consequently, any delayed odd/control datum relevant to the next receiver is not part of the
  current theorem-visible reduced receiver data.
\end{enumerate}
In particular, residual odd completion/source data are \emph{not} part of the current
theorem-visible reduced receiver data.
\end{proposition}

\begin{proof}
Item \textnormal{(i)} is exactly item \textnormal{(R1)} of
\cref*{red-thm:accepted-record-reduction}.
Item \textnormal{(ii)} is item \textnormal{(R4)} of \cref*{red-thm:accepted-record-reduction}.
Then item \textnormal{(iii)} follows: once the current theorem-visible reduced receiver object is
fixed, any odd/control datum relevant to the next receiver can only enter as delayed downstream data
rather than as part of the current reduced receiver semantics.
\end{proof}

\begin{remark}[Current versus delayed visibility]
\label{rem:delay-zigzag}
By \cref{prop:delay-current-closure}, the current accepted-record history determines the current
theorem-visible reduced receiver object. Residual odd completion/source data therefore first appear,
at the theorem-visible level considered here, through the delayed next-receiver readout.
\end{remark}

\section{Delayed Quotient and Selected Line}

\begin{definition}[Delayed odd-control quotient]
\label{def:delayed-odd-class}
Fix an accepted tick \(k\) and the next accepted receiver step \(k+1\).
Let
\[
\mathcal S^{\mathrm{odd}}_{k\to k+1}
\]
be the real vector space of residual odd completion/source data at tick \(k\) whose next-step
incarnations are admissible welded odd trace-channel sources in the sense of
\cref{stack-def:witness-seam-realization}.
Let
\[
\mathcal T^{\mathrm{inv,cur}}_k
\subset
\mathcal S^{\mathrm{odd}}_{k\to k+1}
\]
be the linear subspace of realization differences invisible to every current reduced receiver datum
in the sense of \cref*{red-def:reduced-datum,red-thm:accepted-record-reduction}.
Define the \emph{delayed odd-control quotient} by
\[
\tau^{\mathrm{comp,delay}}_k
:=
\mathcal S^{\mathrm{odd}}_{k\to k+1}\big/\mathcal T^{\mathrm{inv,cur}}_k.
\]
Elements are written
\(
[\tau]^{\mathrm{delay}}_k
\).
\end{definition}

\begin{definition}[Delayed selected generator and delayed selected line]
\label{def:delayed-selected-line}
Fix an accepted tick \(k\) and the next accepted receiver step \(k+1\).
Assume a detector-factorized witness-indexed seam realization at \(\Lambda_{k+1}\) in the sense of
\cref{stack-def:detector-factorized-seam}.
Define the \emph{delayed selected generator} by the delayed class of the unit welded odd source:
\[
\widehat\ell^{\mathrm{delay}}_k
:=
[\tau^{\mathrm{unit}}_{\Lambda_{k+1}}]^{\mathrm{delay}}_k
\in
\tau^{\mathrm{comp,delay}}_k.
\]
The associated \emph{delayed selected line} is
\[
\widehat L^{\mathrm{sel,delay}}_k
:=
\mathbb R\,\widehat\ell^{\mathrm{delay}}_k
\subset
\tau^{\mathrm{comp,delay}}_k.
\]
\end{definition}

\begin{theorem}[Detector-factorized seam sources generate a canonical delayed selected line]
\label{thm:delayed-selected-line}
Fix an accepted tick \(k\) and the next accepted receiver step \(k+1\).
Assume a detector-factorized witness-indexed seam realization at \(\Lambda_{k+1}\) in the sense of
\cref{stack-def:detector-factorized-seam}.
Then every witness-indexed delayed class lies on the delayed selected line:
\[
\big[\tau^{f,y}_{\Lambda_{k+1}}\big]^{\mathrm{delay}}_k
=
\sigma_{\mathrm{fill}}(f)\,
\overline{\mathrm{obs}}(y)\,
\widehat\ell^{\mathrm{delay}}_k
\qquad
\forall f,\forall y.
\]
In particular, the delayed seam packet is one-dimensional after passing to the delayed quotient,
with canonical generator \(\widehat\ell^{\mathrm{delay}}_k\).
\end{theorem}

\begin{proof}
The detector-factorized seam identity of \cref{stack-def:detector-factorized-seam} gives
\[
\tau^{f,y}_{\Lambda_{k+1}}
=
\sigma_{\mathrm{fill}}(f)\,
\overline{\mathrm{obs}}(y)\,
\tau^{\mathrm{unit}}_{\Lambda_{k+1}}.
\]
Passing to the delayed quotient preserves the same linear relation, hence
\[
\big[\tau^{f,y}_{\Lambda_{k+1}}\big]^{\mathrm{delay}}_k
=
\sigma_{\mathrm{fill}}(f)\,
\overline{\mathrm{obs}}(y)\,
\big[\tau^{\mathrm{unit}}_{\Lambda_{k+1}}\big]^{\mathrm{delay}}_k.
\]
By \cref{def:delayed-selected-line}, the last class is exactly
\(
\widehat\ell^{\mathrm{delay}}_k
\).
\end{proof}

\section{Delayed Realization}

\begin{definition}[Delayed realization map on the selected line]
\label{def:delayed-response-map}
Fix an accepted tick \(k\) and the next accepted receiver step \(k+1\).
Assume the hypotheses of \cref{thm:delayed-selected-line} and, in addition, the hypotheses of
\cref{stack-thm:branchwise-seam-witness} on a connected accepted-branch segment containing
\(\Lambda_{k+1}\).
Define the \emph{delayed realization map} on the delayed selected line by
\[
\mathfrak R^{\mathrm{delay}}_{k\to k+1}
:
\widehat L^{\mathrm{sel,delay}}_k
\longrightarrow
\mathbb R\,u^{\mathrm{LoG}}_{\Lambda_{k+1}},
\qquad
\mathfrak R^{\mathrm{delay}}_{k\to k+1}\!\big(a\,\widehat\ell^{\mathrm{delay}}_k\big)
:=
a\,c_{\Lambda_{k+1}}\,u^{\mathrm{LoG}}_{\Lambda_{k+1}},
\]
where \(c_{\Lambda_{k+1}}\) is the normalized seam coefficient of
\cref{stack-thm:branchwise-seam-witness}.
Equivalently, its reduced welded-image form is
\[
\mathsf W^{\mathrm{delay}}_{k\to k+1}
:
\widehat L^{\mathrm{sel,delay}}_k
\longrightarrow
\mathbb R\,P^{\mathrm{LoG}}_{-,\Lambda_{k+1}},
\qquad
\mathsf W^{\mathrm{delay}}_{k\to k+1}\!\big(a\,\widehat\ell^{\mathrm{delay}}_k\big)
:=
a\,c_{\Lambda_{k+1}}\,P^{\mathrm{LoG}}_{-,\Lambda_{k+1}}.
\]
\end{definition}

\begin{proposition}[Delayed realization formula on witness-indexed representatives]
\label{prop:delayed-realization-formula}
Fix an accepted tick \(k\) and the next accepted receiver step \(k+1\).
Assume:
\begin{enumerate}[leftmargin=2em,label=(\alph*)]
\item the hypotheses of \cref{thm:delayed-selected-line};
\item the hypotheses of \cref{stack-thm:branchwise-seam-witness} on a connected
  accepted-branch segment containing \(\Lambda_{k+1}\);
\item \cref{stack-prop:reduced-welded-ambiguity-trivial} at the next receiver step.
\end{enumerate}
Then for every filler tag \(f\) and witness \(y\),
\[
\mathsf W^{\mathrm{red}}_{\Lambda_{k+1}}\!\big(J^{f,y}_{\Lambda_{k+1}}\big)
=
\sigma_{\mathrm{fill}}(f)\,
\overline{\mathrm{obs}}(y)\,
c_{\Lambda_{k+1}}\,
P^{\mathrm{LoG}}_{-,\Lambda_{k+1}}.
\]
Equivalently,
\[
\mathfrak R^{\mathrm{delay}}_{k\to k+1}\!\big(
\big[\tau^{f,y}_{\Lambda_{k+1}}\big]^{\mathrm{delay}}_k
\big)
=
\sigma_{\mathrm{fill}}(f)\,
\overline{\mathrm{obs}}(y)\,
c_{\Lambda_{k+1}}\,
u^{\mathrm{LoG}}_{\Lambda_{k+1}}.
\]
In particular, the delayed realization depends only on the scalar coefficient of the delayed class
on the selected line and is therefore well defined.
\end{proposition}

\begin{proof}
By item \textnormal{(1)} of \cref{stack-thm:branchwise-seam-witness},
\[
\rho^{\mathrm{dom}}_{\Lambda_{k+1}}\!\big(J^{f,y}_{\Lambda_{k+1}}\big)
=
\sigma_{\mathrm{fill}}(f)\,
c_{\Lambda_{k+1}}\,
\overline{\mathrm{obs}}(y).
\]
By \cref{stack-prop:reduced-welded-ambiguity-trivial},
\[
\mathsf W^{\mathrm{red}}_{\Lambda_{k+1}}(F)
=
\rho^{\mathrm{dom}}_{\Lambda_{k+1}}(F)\,P^{\mathrm{LoG}}_{-,\Lambda_{k+1}}.
\]
Applying this to \(F=J^{f,y}_{\Lambda_{k+1}}\) gives the first displayed formula.
The line-valued formula is the same identity written using the canonical generator
\(
u^{\mathrm{LoG}}_{\Lambda_{k+1}}
\)
of the one-dimensional compressed welded channel given by
\cref{stack-thm:finite-cutoff-canonical-odd-spectral-selection}.

Finally, by \cref{thm:delayed-selected-line},
\[
\big[\tau^{f,y}_{\Lambda_{k+1}}\big]^{\mathrm{delay}}_k
=
\sigma_{\mathrm{fill}}(f)\,
\overline{\mathrm{obs}}(y)\,
\widehat\ell^{\mathrm{delay}}_k,
\]
so the delayed realization depends only on the scalar multiplying the selected generator.
\end{proof}

\begin{remark}[Why the codomain is next-receiver and reduced]
\label{rem:delay-codomain}
By \cref{stack-thm:finite-cutoff-canonical-odd-spectral-selection}, the compressed welded LoG
channel at \(\Lambda_{k+1}\) is one-dimensional with canonical generator
\(u^{\mathrm{LoG}}_{\Lambda_{k+1}}\) and canonical projector
\(P^{\mathrm{LoG}}_{-,\Lambda_{k+1}}\).
By \cref{stack-prop:reduced-welded-ambiguity-trivial}, the next-step welded readout factors through
the reduced welded image.
Hence \cref{def:delayed-response-map} lands naturally in the next-step reduced welded LoG line.
\end{remark}

\section{Selected-Line Delayed Sufficiency}

\begin{definition}[Theorem-trivial selected delayed classes]
\label{def:delayed-selected-trivial}
Fix an accepted tick \(k\) and the next accepted receiver step \(k+1\).
Assume the hypotheses of \cref{prop:delayed-realization-formula}.
The \emph{theorem-trivial selected delayed classes} are
\[
\mathcal T^{\mathrm{inv,delay,sel}}_k
:=
\ker\!\bigl(
\mathsf W^{\mathrm{delay}}_{k\to k+1}
\!\mid_{\widehat L^{\mathrm{sel,delay}}_k}
\bigr).
\]
Equivalently, because the current theorem-visible reduced receiver content has already been quotiented
out in \(\tau^{\mathrm{comp,delay}}_k\), these are exactly the selected delayed classes whose
next-receiver reduced welded image vanishes.
\end{definition}

\begin{theorem}[Selected-line delayed response sufficiency]
\label{thm:delayed-response-sufficiency}
Fix an accepted tick \(k\) and the next accepted receiver step \(k+1\).
Assume the hypotheses of \cref{prop:delayed-realization-formula} and
\[
c_{\Lambda_{k+1}}\neq 0.
\]
Then on the delayed selected line,
\[
\ker\!\bigl(
\mathfrak R^{\mathrm{delay}}_{k\to k+1}
\!\mid_{\widehat L^{\mathrm{sel,delay}}_k}
\bigr)
=
\mathcal T^{\mathrm{inv,delay,sel}}_k
=
\{0\}.
\]
\end{theorem}

\begin{proof}
By \cref{def:delayed-response-map}, the maps
\(
\mathfrak R^{\mathrm{delay}}_{k\to k+1}
\)
and
\(
\mathsf W^{\mathrm{delay}}_{k\to k+1}
\)
are the same scalar map written in the canonical line generators
\(
u^{\mathrm{LoG}}_{\Lambda_{k+1}}
\)
and
\(
P^{\mathrm{LoG}}_{-,\Lambda_{k+1}}
\),
respectively.
Hence they have the same kernel on
\(
\widehat L^{\mathrm{sel,delay}}_k
\).
By \cref{def:delayed-selected-trivial}, that common kernel is exactly
\(
\mathcal T^{\mathrm{inv,delay,sel}}_k
\),
which proves the first equality.

By \cref{def:delayed-response-map}
\[
\mathfrak R^{\mathrm{delay}}_{k\to k+1}\!\big(
a\,\widehat\ell^{\mathrm{delay}}_k
\big)
=
a\,c_{\Lambda_{k+1}}\,u^{\mathrm{LoG}}_{\Lambda_{k+1}}.
\]
Since
\(
u^{\mathrm{LoG}}_{\Lambda_{k+1}}\neq 0
\)
by \cref{stack-thm:finite-cutoff-canonical-odd-spectral-selection},
the image vanishes if and only if \(a=0\).
So the kernel is trivial, and therefore the theorem-trivial selected delayed classes are trivial as
well.
\end{proof}

\begin{remark}[Nondegeneracy clause in the delayed selected-line package]
The delayed selected-line package used downstream consists of
\cref{def:delayed-selected-line,def:delayed-response-map,thm:delayed-response-sufficiency}.
The nonvanishing clause \(c_{\Lambda_{k+1}}\neq 0\) is therefore carried as an explicit package
hypothesis at the delayed selected-line stage.
Its source remains the upstream seam-overlap/transport lane rather than a proof internal to the
present note.
\end{remark}

\begin{remark}[Scope of delayed selected-line sufficiency]
\label{rem:delay-vs-same-tick}
Theorem~\cref{thm:delayed-response-sufficiency} is a statement on the selected delayed line
\(\widehat L^{\mathrm{sel,delay}}_k\).
It does \emph{not} identify theorem-triviality on the whole delayed odd-control quotient
\(\tau^{\mathrm{comp,delay}}_k\), and it does \emph{not} replace the intrinsic same-tick
assumption \cref{stack-ass:response-sufficiency} inside the strong characterization package.
\end{remark}

\begin{remark}[Forward bridge boundary]
\label{rem:delay-forward-bridge}
The present note stops at faithfulness on the delayed carrier
\(
\widehat L^{\mathrm{sel,delay}}_k
\)
itself.
Passing from that delayed selected line to a descendant point in
\(
\mathfrak D^-_{\Lambda_{k+1}}
\)
is the content of the later split-controlled scalar bridge
\cref{splitlift-thm:splitlift-scalar-descendant-lift}.
Any further selected quotient or mod-null reentry lies still beyond the scope of the present note.
\end{remark}

\end{document}
